% This document is compiled using pdfLaTeX
% You can switch XeLaTeX/pdfLaTeX/LaTeX/LuaLaTeX in Settings

\documentclass{article}
\usepackage{ctex}
\usepackage{amsmath} % 确保引入了此宏包以支持 aligned

\title{因子日历2026}

\date{\today}

\begin{document}

\maketitle

\section{筹码集中度(行为金融因子-筹码分布)}

\subsection{因子定义}
\begin{equation}
    \text{方式1: } \mathrm{chip\_distri} = \frac{\mathrm{price95}_t - \mathrm{price05}_t}{\mathrm{Chip\_Highest}}
\end{equation}
\begin{equation}
    \text{方式2: } \mathrm{chip\_distri} = \frac{2 \times (\mathrm{price95}_t - \mathrm{price05}_t)}{\mathrm{price95}_t + \mathrm{price05}_t}
\end{equation}

\subsection{变量定义}
\begin{itemize}
    \item ① 假设某只股票在 \( t \) 日的筹码分布共有 \( n \) 个价位（\( i = 1, 2, \ldots, n \)），\( \mathrm{vol}_{i,t} \) 为第 \( t \) 日 \( i \) 价位上的筹码峰高度，\( \mathrm{price}_{i,t} \) 为该筹码峰对应的价格，\( \mathrm{Chip\_Highest}_t \)、\( \mathrm{Chip\_Lowest}_t \) 分别对应筹码最高价和筹码最低价，\( \mathrm{price95}_t \)、\( \mathrm{price05}_t \) 分别为筹码分布 95\% 和 5\% 分位数价格；
    \item ② 上述计算的是 90\% 筹码集中度因子；此外，还可利用筹码分布 85\% 和 15\% 分位数价格计算 70\% 的筹码集中度因子；利用筹码分布 75\% 和 25\% 分位数价格计算 50\% 的筹码集中度因子等。
\end{itemize}

\subsection{说明}
筹码集中度可用于追踪股票吸筹、拉升、出货和下跌的阶段表现；吸筹阶段表现为低位集中筹码建立底仓，但该阶段通常比较隐蔽，换手率不会显著放大，筹码集中度变化不大；拉升/下跌阶段表现为股价脱离筹码集中区域，向上/向下打开筹码价格区间，使集中的筹码向上/向下转移发散，表现为筹码集中度降低；出货阶段对应的是筹码高位充分换手，某一价格区间的成交量放大，换手率显著放大，筹码集中度显著提升。因子值越小，筹码集中度越高，越有可能处于出货阶段，未来表现可能不佳。

\subsection{参考文献}
\begin{enumerate}
    \item 陈升锐, 姚紫薇. 2025. 筹码分布因子系统构建. 中信出版社.
\end{enumerate}

\section{谷岭加权价格比(高频因子-量价相关性类)}

\subsection{因子定义}


1. 对过去 20 日同时点成交量按照 1 倍标准差的标准划分，高于 1 倍标准差划分为“喷发成交量”，低于 1 倍标准差划分为“温和成交量”；

2. 对“喷发成交量”进一步按照连续与否划分：若当前分钟为“喷发成交量”，且前后 1 分钟均为“温和成交量”，则划为“孤立喷发成交量”，称为“量峰”；若当前分钟为“喷发成交量”，且前后 1 分钟也存在“喷发成交量”，则划为“连续喷发成交量”，称为“量岭”；将“温和成交量”称为“量谷”。

3. 计算如下因子的 20 日价格比均值即为谷岭加权价格比因子：
\begin{equation}
    \text{谷岭加权价格比} = \frac{\text{每日量谷成交量加权价格}}{\text{同日量岭成交量加权价格}}
\end{equation}

\subsection{变量定义}
\begin{itemize}
    \item ① 类似的，也可将每日量峰成交量加权价格与量岭成交量加权价格做比较，得到峰岭加权价格比因子，该因子也是正向因子。
\end{itemize}

\subsection{说明}
量谷时点，交易低迷，价格过度反应概率较低；量岭时点，个人投资者交易使得价格过度反应概率较高；因子值越高，个人投资者过度反应使得价格相对交易低迷时期价格向下偏离程度越高，个股未来表现越好。

\subsection{参考文献}
\begin{enumerate}
    \item 魏建榕, 王志豪. 2025. 高频成交量的峰、岭、谷信息. 开源证券.
\end{enumerate}

\section{毒流动性(高频因子-流动性类)}

\subsection{因子定义}

基于逐笔委托中的撤单数据，计算如下毒流动性因子：
\begin{equation}
    \text{毒流动性} = \frac{5\text{秒内撤单数量}}{30\text{秒内撤单数量}}
\end{equation}

\subsection{变量定义}
\begin{itemize}
    \item ① 可对计算的因子计算20日均值进行低频化处理。
\end{itemize}

\subsection{说明}
毒流动性因子从“高频撤单”角度来追踪机构交易行为（散户投资者通常不容易接触到程序化交易和高速交易通道，在订单簿中挂单时间越短的撤单，同样是机构行为的概率也越大）；高频撤单的相对比例越高，说明在盘口集中的撤单中，部分机构在股票中更多地起到了做市商的作用，然而这些流动性并非是可以获得的简单筹码，当趋势呈现反向时他们会迅速地改变原有交易方向，因此这些流动性也是“有毒的流动性”，高频流动性越多，实则股票交易的流动性越差。

\subsection{参考文献}
\begin{enumerate}
    \item 魏建榕, 苏良. 2024. 订单流系列：撤单行为规律初探. 开源证券.
\end{enumerate}

\section{相似低波(量价因子改进)}

\subsection{因子定义}
① 滚动计算历史收盘序列（长度为 120 个交易日）中任意时刻 \( T \) 起连续 \( \mathrm{RW}=6 \) 日序列同当前时刻 \( t \) 的股票收盘价序列的皮尔逊相关系数：
\begin{equation}
    \mathrm{Correlation}_{i,T} = \rho_{\mathrm{Sequence}_{i,t}, \mathrm{TempHistory}_{i,T}}
\end{equation}
\begin{equation}
    \mathrm{Sequence}_{i},t = [\mathrm{Close}_{i,t-\mathrm{RW}+1}, \mathrm{Close}_{i,t-\mathrm{RW}+2}, \dots, \mathrm{Close}_{i,t}]
\end{equation}
\begin{equation}
    \mathrm{TempHistory}_{i,T} = [\mathrm{Close}_{i,T}, \mathrm{Close}_{i,T+1}, \dots, \mathrm{Close}_{i,T+\mathrm{RW}-1}]
\end{equation}

② 设定相关系数的阈值 \( \mathrm{Threshold}=0.4 \) 筛选出走势相似度较高的时刻：
\begin{equation}
    \{ \tau \mid |\mathrm{Correlation}_{i,\tau}| \geq \mathrm{Threshold} \}
\end{equation}

③ 计算历史出现相似走势后既定持仓时间内股票累计超额收益率标准差的倒数即为相似低波因子：
\begin{equation}
    \mathrm{SimilarLowVolatility}_{i,t} = \frac{1}{\mathrm{std}(\mathrm{ER}_{i,\tau})}
\end{equation}
\begin{equation}
    \mathrm{ER}_{i,\tau} = \prod_{k=\tau+\mathrm{RW}}^{\tau+\mathrm{RW}+\mathrm{HoldingTime}-1} (1 + \mathrm{ER}_{i,k}) - 1, \quad 1 < \mathrm{HoldingTime} < \mathrm{RW}
\end{equation}
\begin{equation}
    \omega_{i,k} = \frac{e^{-\lambda k}}{\sum_{k=1}^{n} e^{-\lambda k}}, \quad \lambda = \frac{\ln(2)}{H}
\end{equation}

\subsection{说明}
同一只股票的历史收盘价序列中往往会出现多个走势相似的时刻，每次出现相似走势后，可以计算得到出现相似走势后股票一段时间内的累计超额收益率，多次相似走势后得到的累计超额收益率样本的波动率越低，未来继续持有该股票获得的超额收益率的期望越高。

\subsection{参考文献}
\begin{enumerate}
    \item 郑琳琳, 王天业. 2024. 基于历史相似走势的因子选股研究. 西南证券.
\end{enumerate}

\section{大单推动涨幅(高频因子-动量反转类)}

\subsection{因子定义}


\begin{equation}
    \prod_{n=t}^{t-T+1} \left( prod\left( 1 + r_{i,j,n} \cdot I_{\{j\in IdxSet\}}\right) \right) - 1
\end{equation}

\begin{equation}
    \text{平均单笔成交金额} = \frac{\sum_{j=1}^{N} \mathrm{Amt}_{i,j,n}}{\sum_{j=1}^{N} \mathrm{TrdNum}_{i,j,n}}
\end{equation}

\subsection{变量定义}
\begin{itemize}
    \item ① \( r_{i,j,n} \) 为股票 \( i \) 在第 \( n \) 个交易日内第 \( j \) 分钟的收益率序列；
    \item ② \( \mathrm{Amt}_{i,j,n} \) 为股票 \( i \) 在第 \( n \) 个交易日内第 \( j \) 分钟的成交额序列；
    \item ③ \( \mathrm{TrdNum}_{i,j,n} \) 为股票 \( i \) 在第 \( n \) 个交易日内第 \( j \) 分钟的成交笔数序列；
    \item ④ \( \mathrm{IdxSet} \) 为第 \( n \) 个交易日内平均单笔成交金额最大的 30\% 的分钟 \( K \) 线的序号；
    \item ⑤ 月度选股下，\( T=20 \) 个交易日；周度选股下，\( T=5 \) 个交易日。
\end{itemize}

\subsection{说明}
平均单笔成交金额较大的 \( K \) 线多空博弈激烈，未来的反转效应更强；该因子与股票未来收益负相关。

\subsection{参考文献}
\begin{enumerate}
    \item 冯佳睿，姚石.2019.日内分时交易中的玄机.海通证券
\end{enumerate}

\section{跌幅最大时刻-主卖笔均成交金额（高频因子-成交分布类）}

\subsection{因子定义}


1. 将日内 240 个 1 分钟数据集合标记为 \( T = \{t_1, t_2, \dots, t_{240}\} \)，判断股票 i 的每分钟涨跌幅，将跌幅最大的 10\% 时刻标记为 \( P = \{t_{i1}, t_{i2}, \dots, t_{ik}\} \)；

2. 利用逐笔成交数据，进一步对 P 中时刻的所有成交记录划分为主动买入成交 \( P\_Buy \) 和主动卖出成交 \( P\_Sell \) 两类；

3. 对跌幅最大时刻下的主动卖出成交额和成交笔数进行汇总，根据其成交额 \( \mathrm{Amt}_t \) 之和除以成交笔数 \( \mathrm{DealNum}_t \) 之和，构造跌幅最大时刻且主卖的笔均成交金额 \( \mathrm{ATD}_{\mathrm{DownTop10\%\_Sell}} \)：
\begin{equation}
    \mathrm{ATD}_{\mathrm{DownTop10\%\_Sell}} = \frac{\sum_{t \in P\_\mathrm{Sell}} \mathrm{Amt}_t}{\sum_{t \in P\_\mathrm{Sell}} \mathrm{DealNum}_t}
\end{equation}

4. 同理，将全天的成交金额除以全天的成交笔数，得到全天的笔均成交金额 \( \mathrm{ATD}_T \)：
\begin{equation}
    \mathrm{ATD}_T = \frac{\sum_{t \in T} \mathrm{Amt}_t}{\sum_{t \in T} \mathrm{DealNum}_t}
\end{equation}

5. 将跌幅最大时刻且主卖的笔均成交金额除以全天的笔均成交金额，即得到去量纲化后的标准化因子 \( \mathrm{SATD}_{\mathrm{DownTop10\%\_Sell}} \)：
\begin{equation}
    \mathrm{SATD}_{\mathrm{DownTop10\%\_Sell}} = \frac{\mathrm{ATD}_{\mathrm{DownTop10\%\_Sell}}}{\mathrm{ATD}_T}
\end{equation}

\subsection{说明}
笔均成交额类因子通过汇总“特殊的、具有较多信息含量”时刻的成交金额情况来刻画主力资金的行为动向；笔均成交金额越大，说明该特殊时间内资金优势越明显，其属于主力资金的可能性越大。“跌幅最大时刻”通过涨跌幅度来追踪主力资金抄底意愿强度，叠加逐笔订单主卖意愿做更细切分能进一步提升因子表现（卖方主导下，跌幅越大，主力抄底意愿越强，未来收益越高）。

\subsection{参考文献}
\begin{enumerate}
    \item 张欣悦, 张宇. 2025. 日内特殊时刻蕴含的主力资金 Alpha 信息. 国信证券.
\end{enumerate}

\section{改进大单交易占比（高级因子-资金流类）}

\subsection{因子定义}
\begin{equation}
    \text{改进大单交易占比} = \frac{(-1) \times \text{大买} \& \text{非大卖} + (-1) \times \text{非大买}\&\text{大卖} + \text{大买}\&\text{大卖}}{\text{总成交量}}
\end{equation}

\subsection{变量定义}
\begin{itemize}
    \item ① 大单和非大单的判断标准为：先将同一委买 ID（或委卖 ID）对应的实际成交量进行加总，再将全部实际成交的委买 ID（或委卖 ID）成交量进行降序排列，最后将成交量大于前 10\% 分位点的委买单（或委卖单）记为“大买单”（或“大卖单”），要剔除开盘集合竞价的成交记录。
\end{itemize}

\subsection{说明}
经测试发现对于“委买单为大单、委卖单为非大单”和“委买单为非大单、委卖单为大单”两种情况下的交易占比因子与未来股价负相关，只有“委买单及委卖单均为大单”的交易占比因子才与未来股价正相关，所以在计算大单交易占比因子时，对前面两种情况下做方向上的调整（通过乘 -1 将对未来股价的负向影响转为正向影响）有助于提升原始大单交易占比因子的表现。

\subsection{参考文献}
\begin{enumerate}
    \item 张欣鹏, 张宇. 2024. 高频订单成交数据蕴含的 Alpha 信息. 国信证券.
\end{enumerate}

\section{时间重心偏离（高频因子-收益分布类）}

\subsection{因子定义}
将涨、跌幅的时间重心单独剥离干扰因子（截面回归中心化处理）：
\begin{equation}
    G_u = f_u(\bar{R}_u, R_1, R_2, R_{\mathrm{overnight}}) + \varepsilon_u
\end{equation}
\begin{equation}
    G_d = f_d(\bar{R}_d, R_1, R_2, R_{\mathrm{overnight}}) + \varepsilon_d
\end{equation}

通过横截面回归，构造“时间差”指标，并取 \( N=20 \) 日均值作为因子：
\begin{equation}
    \varepsilon_d = \alpha + \beta \varepsilon_u + \varepsilon
\end{equation}
\begin{equation}
    \mathrm{TGD} = \frac{1}{N} \sum_{i=1}^N \varepsilon
\end{equation}

\subsection{变量定义}
\begin{itemize}
    \item ① \( G_u \) 和 \( G_d \) 分别为涨幅时间重心和跌幅时间重心，具体计算方式见相关日历页；
    \item ② \( \bar{R}_u \) 和 \( \bar{R}_d \) 分别表示涨、跌幅较大的 17 根分钟 Bar 的平均涨幅和平均跌幅；
    \item ③ \( R_1, R_2, R_{\mathrm{overnight}} \) 分别为 09:31-10:00 区间涨跌幅、10:01-10:30 区间涨跌幅、隔夜涨跌幅。
\end{itemize}

\subsection{说明}
涨、跌时间重心通过分钟涨跌幅对修正的时间戳加权平均得到，反映了股价在日内交易中上涨和下跌的位置，捕捉了股票的交易行为特征，两者的相对位置是一个有效的 Alpha 因子，与未来收益正相关；剥离收益率结构和极端收益反转效应等干扰因子，可以增强因子表现。

\subsection{参考文献}
\begin{enumerate}
    \item 魏建榕, 苏良, 张少成. 2022. 日内分钟收益率的时序特征：逻辑讨论与因子增强. 开源证券.
\end{enumerate}

\section{峰岭成交比（高频因子-成交分布类）}

\subsection{因子定义}

1. 对过去 20 日同时点成交量按照 1 倍标准差的标准划分，高于 1 倍标准差划分为“喷发成交量”，低于 1 倍标准差划分为“温和成交量”；

2. 对“喷发成交量”进一步按照连续与否划分：若当前分钟为“喷发成交量”，且前后 1 分钟均为“温和成交量”，则划为“孤立喷发成交量”，称为“量峰”；若当前分钟为“喷发成交量”，且前后 1 分钟也存在“喷发成交量”，则划为“连续喷发成交量”，称为“量岭”。将“温和成交量”称为“量谷”。

3. 计算 20 日量峰总成交额与量谷总成交额，两者做比即为峰岭成交比因子：
\begin{equation}
    \text{峰岭成交比} = \frac{\text{20日量峰总成交额}}{\text{20日量谷总成交额}}
\end{equation}

\subsection{说明}
峰岭成交比因子衡量了知情交易相对个人投资者交易的相对参与程度，因子值越高，未来收益越高。

\subsection{参考文献}
\begin{enumerate}
    \item 魏建榕, 王志豪. 2025. 高频成交量的峰、谷信息. 开源证券.
\end{enumerate}

\section{小买单主动成交度（高频因子-资金流类）}

\subsection{因子定义}
\begin{equation}
    \frac{1}{T} \sum_{n=t}^{t-T+1} \frac{\sum_{j=1}^{N} \text{小买单主动成交额}_{i,j,n}}{\sum_{j=1}^{N} \text{小买单成交额}_{i,j,n}}
\end{equation}

\subsection{变量定义}
\begin{itemize}
    \item ① 基于逐笔成交数据中的叫买与叫卖单号将逐笔成交数据合成为买卖单数据；
    \item ② 根据买卖单分布界定大小单，将小于多日买卖单成交单数据对数调整后的均值作为小单筛选阈值；
    \item ③ \( i, j, n \) 分别表示第 \( i \) 只股票在第 \( n \) 个交易日内的第 \( j \) 分钟的成交额；
    \item ④ 月度选股下，\( T=20 \) 个交易日；
    \item ⑤ 可将买单替换成卖单，即可计算得到小卖单主动成交度。
\end{itemize}

\subsection{说明}
主动成交度类因子刻画了逐笔买卖单中主动成交的程度；可以对大中小的买卖单都计算成交度，其中小单的相关因子表现最优，无论是小买单还是小卖单，主动成交度越高，股票未来收益表现更优。

\subsection{参考文献}
\begin{enumerate}
    \item 冯佳睿, 袁林青. 2021. 买卖单主动成交中的隐藏信息. 海通证券.
\end{enumerate}
\end{document}